\documentclass[11pt]{article}
\usepackage[utf8]{inputenc}
\usepackage[T1]{fontenc}
\usepackage{geometry}
\usepackage{enumitem}
\usepackage{xcolor}
\usepackage{hyperref}
\usepackage{amsmath}

\geometry{margin=1in}

\newcommand{\tagg}[1]{\textcolor{blue}{[#1]}}

\title{\textbf{Singular Learning Theory for Applied Deep Learning}\\[8pt]
\large Analysis Plan by Section}
\author{Working Document, aligned with the goals document (v5)}
\date{\today}

\begin{document}
\maketitle

\noindent\emph{This plan implements the goals, structure, and framing
decisions in the companion goals document (v5). Each section expands the
structure overview into argument-level detail. Open items are flagged
inline and collected in the needs document (v3); items marked [VERIFY]
carry a lineage or factual claim to confirm before prose; items backing
a goals-document contingency carry its tag [C1]--[C7]. Settled strategy
is stated in the body. Cross-references to goals use aims (i)--(iv) and
walk-away items 1--7 of v5.}

% ============================================================================
\section{Section 1: Something Is Wrong}
% ============================================================================

\textbf{Goal:} Urgency, and only urgency --- and only level-1 urgency:
tools that make parameter-space claims and are wrong at their own stated
jobs. Serves aim (i) level 1 (walk-away items 1--3); plants aim (i)
level 2 as a question (what did my optimizer choose, and from what?);
plants the lens (the reader cannot yet tell geometry from sample size
from computation).

\tagg{DECISION: \S1 is built on four escalating pillars, each a larger
frame than the last: wrong results $\to$ misaimed fixes $\to$
regularizing instead of redesigning $\to$ folklore in the gap. The
ordering is the argument. Evidence calibrated per pillar (demonstrated /
demonstrable / argued-from-pattern / argued-from-vetted-instances). All
four pillars gated to level 1: every exhibit makes a parameter-space
claim.}

\tagg{DECISION: The masking thesis is stated early and carried by every
pillar, with its honest history (goals decision (d)): the repairs were
aimed at stability and delivered it; the masking is a byproduct of a
success, which is why no one noticed and no one is accused. Each repair
returns a plausible number, so the failure is silent. The two visible
symptoms are (a) structurally wrong answers where checkable [C1], (b)
unattributable success elsewhere. Right-tool-wrong-story appears only as
symptom (b), not as a thesis.}

\subsection*{\S1.1: Pillar 1 --- wrong results (the anchor)}

Regular-geometry tools produce demonstrably false outputs on the models
practitioners use. Open here because it is undeniable.

\begin{itemize}[nosep]
\item \textbf{Variational inference.} The Gaussian/mean-field family
gives wrong posteriors on singular models; the Wei--Lau
resolution-of-singularities family is the demonstrated fix. Not
suboptimal --- false, with the fix already shown. The claim is a
parameter-space shape claim: the family asserts the posterior over
weights is regular; the geometry falsifies the assertion.
\item \textbf{Normalization.} On scale-invariant nets
($f(x,\alpha\theta)=f(x,\theta)$) the textbook weight-decay
complexity-reduction rationale is void on the architectures it runs on.
\end{itemize}

\tagg{NEED: exact citations and precise statements; the scale check on
singular VI decides \S4.3's register [C2]. Wei--Lau (arXiv:2302.06035);
van Laarhoven lineage for normalization. Needs doc \S1.}

\tagg{FIGURE: VI-Gaussian reporting a wrong posterior on the microscope
FFN, corrected by the singularity-aware family.}

\subsection*{\S1.2: Pillar 2 --- misaimed fixes}

When a tool misbehaves, the reach is for another regular-geometry patch
aimed at the symptom.

\begin{itemize}[nosep]
\item Influence-function and Laplace damping --- fabricating curvature on
a degenerate Hessian. The Koh--Liang lineage is the cleanest exhibit:
assumption stated (strict convexity, PD Hessian), failure acknowledged,
damping $H+\lambda I$ shipped, and the patch's own failure documented
(the $\lambda$ needed on large ReLU nets distorts the estimate).
\item Sharpness/flatness-seeking --- chasing a
reparameterization-non-invariant metric (Dinh et al.\ 2017); the field's
response repaired the sharpness measure rather than asking what invariant
it shadows.
\item Isotropizing preconditioners / second-order refinements that damp
the degenerate structure.
\end{itemize}

\tagg{DECISION (updated to goals v5): The pillar-tool audit table lands
here (drafted: dl\_section1\_audit\_table.tex). Columns: tool, stated
assumption, shipped repair, documented strain, diagnosis tag
(geometry / sample size / computation). \emph{Three} gates per row:
assumption stated or structurally required in the pillar paper; repair
restores regularity by hand; the row makes a parameter-space claim
(posterior, uncertainty region, influence, curvature, parameter count).
The convergence-guarantees row is \emph{removed from the table} --- it
is a level-2 item and lives in \S4.2's recast (goals decision (k)).}

\tagg{NEED: reconcile the drafted table to this framing --- remove the
convergence row, add the third gate and the diagnosis-tag column. Needs
doc \S1.}

\subsection*{\S1.3: Pillar 3 --- regularizing instead of redesigning}

Not a single wrong fix but the sustained reflex: weight decay, damping,
sharpness penalties are the same move at different sites. Argued from
the pattern the audit table makes visible --- the ``shipped repair''
column is this pillar rendered as data. [C5: the one-reflex reading
must hold across the (now level-1-gated) shipped-repair column;
fallback ``a family of regularity-restoring moves.'']

\subsection*{\S1.4: Pillar 4 --- folklore in the gap}

Lacking a name for the degeneracy, the field routes around it; lore
accumulates where theory should be. Framed structurally, not
accusatorially, and with the accurate history: the structure was seen
(information geometry: Fisher degeneracy on the sets of
unidentifiability, plateaus as its dynamical shadow --- Fukumizu 1996,
Fukumizu--Amari 2000, Amari--Park--Ozeki 2002/2006; Hagiwara noticing
AIC fails on perceptrons), a partial language existed, and the
resolution had to be imported from algebraic geometry (Watanabe).
Folklore is what accumulated in the transfer gap.

\tagg{DECISION (updated to goals v5, decision (k)): two gates per
folklore face: the phenomenon was genuinely encountered; the workaround
genuinely functions as folklore (tuned by feel, justified post hoc).
\S1's faces are \emph{two}: mean-field's variational gap misattributed;
the damping ritual across Laplace / influence / natural gradient /
K-FAC. The convergence-lore face (tuning lore bridging scoped
guarantees to practice) is level 2 and moves to \S4.2.}

\tagg{DECISION: the ``encounter without the language'' paragraph
(drafted with the audit table) carries the observational literature ---
Sagun spectra, mode connectivity, re-basin --- as this pillar's
evidence, kept out of the table because those works make no regularity
assumption. This is also the walk-away item 1 material: the field met
the phenomena and even mapped the set they trace out; what was missing
is the language for the geometry that governs them.}

\tagg{VERIFY: Hagiwara citation and year; keep the Watanabe--Amari
lineage institutional, not biographical, unless a citation is found.
Needs doc \S1.}

\subsection*{\S1.5: The spared tool and the handoff}

Cross-validation flagged as the tool that works; reason deferred to \S4.
Close: one cause, silently masked; a question the reader cannot yet ask
precisely --- \emph{what did my optimizer choose, and from what?} ---
planting level 2 and \S4's instrument; the reader cannot yet tell
geometry from sample size from computation; why does it all fail at
once? Hand off to \S2.

\tagg{DECISION: \S1 is the source for the companion blog.}

% ============================================================================
\section{Section 2: Communicating the Singular Geometry}
% ============================================================================

\textbf{Goal:} Answer \S1's question by installing the object --- the
correct geometric picture without algebra. Serves aim (ii) (walk-away
item 4) and the recognition half of the lens (walk-away item 5's habit).

\tagg{DECISION: three model roles (the DL ladder): microscope FFN
(\S1--2), high-dimensionality reckoning (\S2.3), workbench (\S4).}

\tagg{DECISION (goals decision (a) note): in this plan ``the set''
abbreviates the pair --- the set and how the loss vanishes along it.
Paper prose must not inherit the abbreviation where the vanishing
structure is doing the work; topic sentences, captions, and closes
attribute consequences to the geometry around the set.}

\subsection*{\S2.1: The two engines}

Continuous symmetry (normalization manufactures exact scale-invariance
before any data) and overparameterization / DGP degeneracy. They
compound; the minimizing set's geometry depends on model and data.

\tagg{NEED: catalogue architectural continuous symmetries with
citations; state DGP degeneracy honestly as a presumption. Needs doc
\S2.}

\subsection*{\S2.2: Normal crossing, not wide valley (the \S2 thesis)}

The most-resisted move. Primary signature: hyperbolic level sets
($z_1z_2=\epsilon$), not parabolic ($z_1^2+z_2^2=\epsilon$). Secondary:
stratum-dependent jumps in flat-dimension. And the grading, taught here
because this is where the consequences live (walk-away item 4): the
crossing is built of strata, and the loss vanishes to different orders
along and across them --- $z_1^2$ and $z_1^4$ share a solution set and
differ in everything the paper delivers.

\tagg{DECISION (new in goals v4/v5): underdetermination stated where the
picture makes it obvious --- every point of the crossing is
loss-equivalent, so ``minimize the loss'' does not finish the job;
something else chooses, and \S4 names the instrument that reads the
choice. One paragraph here, at picture altitude; the full recast is
\S4.2's.}

\tagg{VERIFIED: Watanabe's own framing --- a deep-learning sub-model
corresponds to an algebraic variety in parameter space; singularities
are made normal crossing via Hironaka's resolution; the process is
captured by two birational invariants (RLCT, singular fluctuation).
Cite his 2009 book and expository statements. Resolved in needs doc.}

\subsection*{\S2.3: The descent in visibility}

Plot the crossing on the FFN; calibrate the eigenspectrum probe against
the visible contour; climb in dimension until a generic slice goes flat
--- a reported finding, not an embarrassment. Then the reckoning: in
real models the curse of dimensionality forecloses profile likelihood,
visualized posteriors, contour inspection entirely. The foreclosure is
the argument for probes and, ultimately, for SLT. Intuition to install:
when you find flatness, think ``a crossing whose branches I can't all
see,'' not ``wide valley.''

\tagg{FIGURE: (i) hyperbolic vs.\ parabolic contours; (ii) contour
$\leftrightarrow$ eigenspectrum side-by-side; (iii) the slice going flat
with dimension. Optional candidate (needs doc \S2): level-set curvature
concentrating at stratum intersections as a crossing signature --- a
visualization device only; the prose vocabulary stays volume-based.}

\subsection*{\S2.4: Posterior pinned to the singular set; geometry vs.\
$n$ (promoted)}

Show mass conforming to the crossing and the geometry's visibility
changing with sample size: diffuse $\to$ ridge $\to$ concentrated on the
singular set; null space profiled against $n$ and a degeneracy knob.

\tagg{DECISION (mechanism, from goals v5 \S2 overview): the sample-size
arm is taught through the diffusion mechanism --- the level sets of the
population loss fix how much room finite-$n$ estimates have to wander
from the set; small $n$ leaves room to sit far from the singular
structure, where it is blurred; growing $n$ closes the room, revealing
the set. The figure's specified job: the posterior/estimate cloud
shrinking onto the crossing as $n$ grows. This phrasing (``how much room
the level sets leave'') is the phrase-bank bridge to \S3.4's volume
scaling --- reuse verbatim.}

\tagg{DECISION: the figure must support two readings: at small $n$ a
singular model masquerades as regular (your regular tool ``worked''
because $n$ hid the geometry), and singular effects grow as $n$ grows
(conclusions at this $n$ may not transfer). Design the figure to expose
both axes.}

\tagg{VERIFY [C4]: the $n$-revelation and compounding claim against the
finite-sample / dimension-dependence literature before the strong form
is stated; the evidentiary bar is higher now that the diagnosis leans on
it. If verified, the compounding is restored to the goals NOT bullet and
walk-away item 5. Needs doc \S2.}

\subsection*{\S2.5: Why standard fixes don't work}

Reparameterization cannot remove the crossing (topological);
regularization shifts which branch dominates rather than removing
degeneracy; optimizer diagnostics look healthy while the chain mixes
along the variety. Connects to \S1's masking: healthy-looking numbers
are exactly what masking predicts.

\subsection*{\S2.6: The mechanism generalizes $+$ the fork}

The crossing is not special; symmetry and overparameterization are
universal. One sentence: singular geometry can be transformed into a
tractable standard form. Fork: tools (\S4) or mathematics (\S3).

% ============================================================================
\section{Section 3: The Core Framework}
% ============================================================================

\textbf{Goal:} Formalize \S2. Skippable by the tools-only practitioner;
the on-ramp for the theorist. Serves walk-away item 5 (the framework and
the lens it grounds) and aim (ii)'s formal half.

\subsection*{\S3.1: Regular geometry formalized}

Fisher as metric; regularity $\to$ quadratic KL $\to$ Laplace $\to$
$p/2$. Brief.

\subsection*{\S3.2: Singular geometry formalized}

$W_0=\{\theta: K(\theta)=0\}$ a variety; degenerate leading terms;
Fisher degeneracy; the stratified parameter space (the formal version of
\S2's crossing); conditions (analyticity, common support, proper prior).
Illustrated on the FFN with the \S2 pictures relabeled.

\tagg{DECISION (goals decision (t)): the analyticity caveat stated
plainly here --- the theorems require analytic models; the microscope
FFN uses an analytic activation so every step is carried out in full
rigor; ReLU is the limit of analytic units (softplus,
$\beta\to\infty$) whose exact positive homogeneity strengthens the
symmetry argument; the gap is in the machinery, not the phenomenon.
Modern transformer activations (GELU, SiLU) are analytic, narrowing the
gap for the reader's own models [C6]. Formal ReLU extension flagged as
open in \S5.}

\tagg{VERIFY: Watanabe's exact definitions and conditions; locations in
Watanabe 2009/2013. Needs doc \S3.}

\subsection*{\S3.3: The blow-up}

Resolution: $K(g^{-1}(z))=z_1^{2k_1}\cdots z_d^{2k_d}u(z)$;
before/after contrast (tangled crossing $\to$ rectangular level sets);
Hironaka used, not constructed; $\lambda$ invariant across resolutions.
The sorting principle: integrals survive the non-smooth change of
variables; MLE/Fisher/Wald/Hessian-at-a-point do not. \tagg{DECISION:
state the sorting principle as the formal ground of the geometry arm of
the lens --- a geometry problem is one where the local-differential
object was never the right object.}

\tagg{NEED: explicit blow-up for the FFN crossing; before/after figure.
Needs doc \S3.}

\subsection*{\S3.4: Consequences}

Factorized integral; prior and Jacobian enter together; $\lambda=\min_j
(h_j+1)/(2k_j)$; $F_n=nL_n(\hat\theta)+\lambda\log n-(m-1)\log\log
n+O(1)$; tempered posteriors, WBIC $\beta=1/\log n$. The \S2.4
concentration rate gets its number.

\tagg{DECISION (phrase bank): introduce $\lambda$ through the volume
scaling $\mathrm{Vol}\{K \le \varepsilon\} \sim
c\,\varepsilon^{\lambda}(-\log\varepsilon)^{m-1}$ in the \S2.4
vocabulary: ``the room the level sets leave shrinks like
$\varepsilon^{\lambda}$.'' The walk-away pre-teaches this; the payoff
requires the same words.}

\tagg{VERIFY: theorem statements; $\lambda$ for the FFN (compute). Needs
doc \S3.}

\subsection*{\S3.5: Finite-sample validation}

SLT prediction vs.\ thermodynamic integration vs.\ parameter-counting
baseline at multiple $n$ on the FFN; honesty about the $n$ needed at
realistic scale (a computation datapoint for \S4).

% ============================================================================
\section{Section 4: Now What}
% ============================================================================

\textbf{Goal:} The delivered value. A maturity-aware field report
organized around the questions tools answer (walk-away item 6); each
tool tagged by the lens (geometry / sample size / computation) and
graded on concept- and computation-maturity separately, computability
judged at the reader's scale.

\tagg{NEED: \S4 cannot be written from impressions. Two gating tasks
(goals decision (o)): the tools-and-maturity audit (per tool ---
problem-solved / diagnosis tag / concept-maturity /
computation-maturity / scale / progress / rough edges / reader layer)
and the practice inventory (per activity --- claim made / tool used /
where the computation lives / question answered vs.\ question asked /
verdict). Needs doc \S4.}

\subsection*{\S4.1: Two questions, two criteria (the organizing recast)}

\tagg{DECISION (updated from ``CV vindicated'' to the goals v5 hinge,
decision (h)): the practitioner layer opens on the two-questions recast,
not on CV alone. If the question is \emph{prediction}, CV/WAIC is the
right object and is vindicated in singular models (Watanabe 2010: Bayes
CV $\approx$ WAIC including in singular models; gen$+$CV $\to
2\lambda/n$; AIC/DIC, assuming normality, provably fail) --- CV works
because it is predictive/integral. If the question is \emph{identifying
structure}, the free energy is the right object, WBIC estimates it, and
$\lambda$-based selection is consistent where CV/WAIC is not; Watanabe
states both halves and himself deferred the selection comparison. ``CV
vindicated'' must never slide into ``CV suffices'' --- the spared tool
is spared under a scope. Quantitative small-$n$ statement (selection
signal shrinking against CV noise) [C3]: stated only if pinned;
fallback states the consistency asymmetry without the rate.}

\tagg{VERIFY: exact 2010 statements; the C3 citation hunt (Watanabe 2021
mixture comparison; Yamazaki--Kaji 2013) or the mixture experiment
(C1a). Needs doc \S4.}

\subsection*{\S4.2: Practitioner layer --- recast tools $+$ scope table}

Adam, weight decay, flatness recast. \tagg{DECISION (updated to goals
v5): the Adam recast runs through \emph{underdetermination} (decision
(e)): the guarantee is honest (an $\epsilon$-stationary point under
(generalized) smoothness; guarantees promising minima invoke PL/strong
convexity) --- and the target is not a well-posed destination: the
objective specifies a set; the optimizer and its hyperparameters
silently complete the choice. Whether $\epsilon$ and its kin merely
navigate the degenerate landscape or select the stratum is posed as the
open question it is, with Adam's $\epsilon$ flagged as a candidate
damping instance [VERIFY]. Twin calibration with CV (decision (e)):
Adam's guarantee certifies a stationary point, nothing about which; CV
is a generalization estimator, inconsistent for selection --- the
workhorse and the spared tool submit to the same scope-narrowing move.
The convergence-lore folklore face (tuning lore bridging scoped
guarantees to practice) lands here, moved from \S1.}

\tagg{DECISION: the tool-by-question scope table lands here (what it
answers / what practice asks of it / the gap), with safe/caution/avoid
as its practitioner projection --- built from the practice inventory,
not impressions. The practitioner can stop here.}

\tagg{VERIFY: the Adam $v_t$--degenerate-subspace alignment stays an
illustration unless verified; the $\epsilon$-as-damping flag checked
before stated. Needs doc \S4.}

\subsection*{\S4.3: Methods layer --- how to compute}

\begin{itemize}[nosep]
\item \textbf{The LLC as the readout of where training landed} (leads
the layer; walk-away item 6's ``choice made measurable''): the
SGD-vs-entropy-SGD exhibit (1.9M-parameter FFN on MNIST; 80 repetitions;
entropy-SGD lands on measurably lower-$\lambda$ strata with remarkably
low variance) --- optimizer inductive bias made visible by the number
the loss cannot see; NGD-vs-SGD as the second instance. The emblematic
framing (``equivalent by every number you currently look at; different
by the number you don't'') is [C7]-contingent on matched fit; fallback
reports inductive-bias-made-measurable without the equivalence.
\item \textbf{Singular VI} (Wei--Lau): the \S1 indictment becomes the
tool. Diagnosis tag: geometry (the family was the wrong object); the
register depends on the scale check [C2]: provocative if computable at
the reader's scale, else geometry-settled / computation-frontier and
humble (decision (i)).
\item \textbf{The estimation toolkit}: LLC/rLLC, \texttt{devinterp},
susceptibilities. Diagnosis tag: computation (right object; scale is the
frontier); the humble register.
\item \textbf{Bayesian influence functions}: the twin arc --- the damped
IF of \S1's ritual returns as the leading-order Laplace term of the
well-posed BIF, the folklore $\lambda$ as the principled localization
$\gamma$. \tagg{DECISION: make the Wei--Lau and BIF arcs deliberately
parallel: indictment $\to$ tool, twice.}
\item \textbf{Sampler craft}: tuning temperature, localization, weight
restriction so the posterior explores the strata --- mixing along the
variety, not stuck on one branch, not so fast the geometry is
unresolved. Presented as emerging craft with real tooling; a formally
stratification-preserving sampler is flagged as open (feeds \S5).
\end{itemize}

\tagg{NEED: tool roster and grades from the audit; sampler-tuning depth
is the deepest assessment; the C7 matched-fit check. Needs doc \S4.}

\subsection*{\S4.4: Worked example and close}

Refined-LLC for attention heads (arXiv:2410.02984) as the workbench
demonstration, in the appendix with a one-paragraph pointer; doubles as
the bridge to \S5's interpretability flag. Close: which tools to trust
for which question, which are real-but-early, which kind of hardness
each addresses; every good instinct, including the spared tool, is right
under a scope. The delivered value closes here.

% ============================================================================
\section{Section 5: Conclusions and the Open Horizon}
% ============================================================================

\textbf{Goal:} Conclude, then pose the open question and plant the flag
(aim (iv), walk-away item 7). Unification is a question, not a result.

\subsection*{\S5.1: Recap and loss-landscape consequences}

Phase transitions (regularizers select strata; a mechanism for part of
double descent); quasi-singular phenomena (RLCT growth; bearing on
scaling).

\tagg{VERIFY: phase-transition computation on a DL model (or borrow
conceptually from the binomial mixture); quasi-singular RLCT-growth
status in the literature. Needs doc \S5.}

\subsection*{\S5.2: The specification gap as a research program}

\tagg{DECISION (new in goals v4/v5; previously absent from this plan):
hyperparameters silently complete a choice the objective cannot express;
the LLC makes the choice measurable (\S4.3's exhibit); an optimizer
targeting the free energy already demonstrates explicit selection ---
so principled, even optimal, hyperparameters are a well-posed open
question. Explicitly conjectural; the hyperparameter-dependence of
final $\lambda$ within a single algorithm flagged as an open
experiment, not a paper dependency.}

\tagg{VERIFY: pin the free-energy-targeting optimizer citation and what
it demonstrates vs.\ what the program conjectures. Needs doc \S5.}

\subsection*{\S5.3: The open question and the substantive flag}

Could singular geometry unify deep learning practice? Posed genuinely;
the loss-landscape case makes it credible. The flag (sketch, don't
prove): mechanistic interpretability developed at half-page depth
(non-unique circuits, superposition as degeneracy on the singular set,
the rLLC bridge from \S4.4); activation/representation space
(symmetries breaking manifold-hypothesis assumptions and regular-base
flows) and policy/value space (behavior-preserving reparameterizations;
natural policy gradient inheriting the local-differential failure)
compressed to a credible sentence each.

\tagg{NEED: each sketch credible, not closed; susceptibility references
(Baker et al.\ 2025). Needs doc \S5.}

\subsection*{\S5.4: Art to science}

\S1 showed the folklore; \S5 names the science. Capstone: the principles
had to be imported (Hironaka), and the import is not finished --- the
\S4 sampler frontier points toward geometry-preserving machinery of the
kind physics already has for singular configuration spaces (explicitly
horizon, not claim); the ReLU extension named alongside it as open
mathematics. Close: the loss landscape is governed by singular geometry
and there are tools to use now; the unification is open --- the case is
strong enough to pursue. Invitation, not synthesis.

\end{document}
